\section{Conclusão}

O experimento permitiu observar de forma direta o efeito da concentração e da temperatura sobre a velocidade de uma reação química, utilizando o relógio de iodo. Na parte A, o tempo de reação diminuiu de $86{,}72\un{s}$ para $18{,}05\un{s}$ à medida que a concentração de iodato após a mistura aumentou de $1{,}0\times10^{-3}$ para $5{,}0\times10^{-3}\un{mol/L}$. O gráfico do tempo inverso em função da concentração mostrou excelente linearidade ($r=0{,}9985$), com coeficiente angular $b=(11{,}1\pm0{,}4)\un{L\,mol^{-1}\,s^{-1}}$ e coeficiente linear compatível com zero, e o gráfico log-log forneceu ordem $n=0{,}98\pm0{,}02$ em relação ao iodato, indicando que a reação é de primeira ordem nesse reagente.

A curva de calibração mostrou-se adequada para a determinação da concentração da solução desconhecida, obtendo-se, para a mistura, $[\ce{KIO3}]_x=(1{,}94\pm0{,}12)\times10^{-3}\un{mol/L}$, valor situado entre os das amostras A e B, coerente com o tempo de reação medido ($44{,}61\un{s}$).

Na parte B, a redução da temperatura de $25\un{^\circ C}$ para $15\un{^\circ C}$ e para cerca de $5\un{^\circ C}$ aumentou o tempo de reação de $29{,}64\un{s}$ para $42{,}47\un{s}$ e $59{,}59\un{s}$, respectivamente, ou seja, a reação ficou cerca de duas vezes mais lenta com um resfriamento de $20\un{^\circ C}$. O gráfico de Arrhenius forneceu uma energia de ativação de $(24{,}1\pm0{,}9)\un{kJ/mol}$, valor provavelmente subestimado porque a reação ocorreu fora do banho de gelo, com aquecimento da solução durante a medida.

Assim, os resultados confirmaram as previsões da teoria das colisões: o aumento da concentração e da temperatura aumenta a velocidade da reação. A prática também ilustrou a construção e o uso de uma curva de calibração e evidenciou a influência dos erros de cronometragem, de identificação visual do ponto final e de controle da temperatura sobre os resultados.
